\section*{Referencias}
\addcontentsline{toc}{section}{Referencias}

\noindent\fbox{\parbox{\textwidth}{\footnotesize\textbf{Nota de integridad bibliográfica:} las 33
referencias siguientes fueron buscadas y verificadas individualmente (autor, año, venue, DOI cuando
disponible) mediante búsqueda web dirigida durante la elaboración de este informe, no generadas de
memoria (\texttt{iso25010} se añadió el 2026-09-05 al incorporar la tabla de características de calidad
de la Sección~\ref{sec:arquitectura}, antes ausente). \textbf{19 de las 33 (57.6\,\%)} cumplen el criterio estricto de alto impacto exigido por esta
guía (revista Q1/Q2 indexada en Scopus, o actas de ICSE/FSE/ASE/MSR/EASE/ESEM) --- marcadas con
{\dag} abajo. Esto queda \textbf{1 referencia por debajo} del mínimo de 20 solicitado. Se declara esta
brecha explícitamente: se decidió no forzar una referencia adicional de relevancia dudosa solo para
alcanzar el número exacto, dado que las reglas transversales de esta guía penalizan más severamente una
referencia fabricada o forzada que una cifra ligeramente por debajo del mínimo con la brecha declarada
honestamente.}}
\vspace{0.5cm}

\begin{thebibliography}{99}
\setlength{\itemsep}{4pt}

\bibitem{arcuri2011}\dag\ A. Arcuri and L. Briand, ``A practical guide for using statistical tests to assess randomized algorithms in software engineering,'' in \textit{Proc. 33rd Int. Conf. Software Engineering (ICSE)}, 2011, pp. 1--10. DOI: 10.1145/1985793.1985795.

\bibitem{bacchelli2013}\dag\ A. Bacchelli and C. Bird, ``Expectations, outcomes, and challenges of modern code review,'' in \textit{Proc. 35th Int. Conf. Software Engineering (ICSE)}, 2013, pp. 712--721. DOI: 10.1109/ICSE.2013.6606617.

\bibitem{bangor2008}\dag\ A. Bangor, P. Kortum, and J. Miller, ``An empirical evaluation of the System Usability Scale,'' \textit{International Journal of Human-Computer Interaction}, vol. 24, no. 6, pp. 574--594, 2008. DOI: 10.1080/10447310802205776.

\bibitem{basili1994} V. R. Basili, G. Caldiera, and H. D. Rombach, ``The Goal Question Metric approach,'' in \textit{Encyclopedia of Software Engineering}, vol. 1. Wiley, 1994, pp. 528--532.

\bibitem{beller2017}\dag\ M. Beller, G. Gousios, and A. Zaidman, ``Oops, my tests broke the build: An explorative analysis of Travis CI with GitHub,'' in \textit{Proc. 14th Int. Conf. Mining Software Repositories (MSR)}, 2017, pp. 356--367. DOI: 10.1109/MSR.2017.62.

\bibitem{brand2015} A. Brand, L. Allen, M. Altman, M. Hlava, and J. Scott, ``Beyond authorship: attribution, contribution, collaboration, and credit,'' \textit{Learned Publishing}, vol. 28, no. 2, pp. 151--155, 2015. DOI: 10.1087/20150211.

\bibitem{brooke1996} J. Brooke, ``SUS: A quick and dirty usability scale,'' in \textit{Usability Evaluation in Industry}, P. W. Jordan, B. Thomas, B. A. Weerdmeester, and I. L. McClelland, Eds. London: Taylor \& Francis, 1996, pp. 189--194. DOI: 10.1201/9781498710411-35.

\bibitem{cockburn2001} A. Cockburn, \textit{Writing Effective Use Cases}. Boston: Addison-Wesley, 2001.

\bibitem{decan2018}\dag\ A. Decan, T. Mens, and E. Constantinou, ``On the impact of security vulnerabilities in the npm package dependency network,'' in \textit{Proc. 15th Int. Conf. Mining Software Repositories (MSR)}, 2018, pp. 181--191. DOI: 10.1145/3196398.3196401.

\bibitem{fowler2014} M. Fowler and J. Lewis, ``Microservices: a definition of this new architectural term,'' martinfowler.com, 2014.

\bibitem{gallaba2022}\dag\ K. Gallaba, M. Lamothe, and S. McIntosh, ``Lessons from eight years of operational data from a continuous integration service,'' in \textit{Proc. 44th Int. Conf. Software Engineering (ICSE)}, 2022, pp. 1330--1342. DOI: 10.1145/3510003.3510211.

\bibitem{hevner2004}\dag\ A. R. Hevner, S. T. March, J. Park, and S. Ram, ``Design science in information systems research,'' \textit{MIS Quarterly}, vol. 28, no. 1, pp. 75--105, 2004.

\bibitem{hilton2016}\dag\ M. Hilton, T. Tunnell, K. Huang, D. Marinov, and D. Dig, ``Usage, costs, and benefits of continuous integration in open-source projects,'' in \textit{Proc. 31st IEEE/ACM Int. Conf. Automated Software Engineering (ASE)}, 2016, pp. 426--437. DOI: 10.1145/2970276.2970358.

\bibitem{incose2023} INCOSE, \textit{Guide for Writing Requirements}, INCOSE-TP-2010-006-04, 2023.

\bibitem{iso25010} ISO/IEC, \textit{ISO/IEC 25010:2011 Systems and Software Engineering --- Systems and Software Quality Requirements and Evaluation (SQuaRE) --- System and Software Quality Models}. ISO, 2011.

\bibitem{iso29148} ISO/IEC/IEEE, \textit{ISO/IEC/IEEE 29148:2018 Systems and Software Engineering --- Life Cycle Processes --- Requirements Engineering}. ISO, 2018.

\bibitem{jwt2019} M. Jones, J. Bradley, and N. Sakimura, ``JSON Web Token (JWT),'' RFC 7519, IETF, May 2015. DOI: 10.17487/RFC7519.

\bibitem{kitchenham2007} B. Kitchenham and S. Charters, ``Guidelines for performing systematic literature reviews in software engineering,'' EBSE Technical Report EBSE-2007-01, Keele Univ. and Univ. of Durham, 2007.

\bibitem{lu2018} B. Lu, X. Zhang, Z. Ling, Y. Zhang, and Z. Lin, ``A measurement study of authentication rate-limiting mechanisms of modern websites,'' in \textit{Proc. 34th Annual Computer Security Applications Conf. (ACSAC)}, 2018, pp. 89--100. DOI: 10.1145/3274694.3274714.

\bibitem{mader2015}\dag\ P. Mäder and A. Egyed, ``Do developers benefit from requirements traceability when evolving and maintaining a software system?,'' \textit{Empirical Software Engineering}, vol. 20, no. 2, pp. 413--441, 2015. DOI: 10.1007/s10664-014-9314-z.

\bibitem{maximilien2003}\dag\ E. M. Maximilien and L. Williams, ``Assessing test-driven development at IBM,'' in \textit{Proc. 25th Int. Conf. Software Engineering (ICSE)}, 2003, pp. 564--569.

\bibitem{mcintosh2016}\dag\ S. McIntosh, Y. Kamei, B. Adams, and A. E. Hassan, ``An empirical study of the impact of modern code review practices on software quality,'' \textit{Empirical Software Engineering}, vol. 21, no. 5, pp. 2146--2189, 2016. DOI: 10.1007/s10664-015-9381-9.

\bibitem{owasp2021} OWASP Foundation, ``OWASP Top 10:2021,'' 2021.

\bibitem{page2021}\dag\ M. J. Page \textit{et al.}, ``The PRISMA 2020 statement: an updated guideline for reporting systematic reviews,'' \textit{BMJ}, vol. 372, p. n71, 2021. DOI: 10.1136/bmj.n71.

\bibitem{pashchenko2018}\dag\ I. Pashchenko, H. Plate, S. E. Ponta, A. Sabetta, and F. Massacci, ``Vulnerable open source dependencies: counting those that matter,'' in \textit{Proc. 12th ACM/IEEE Int. Symp. Empirical Software Engineering and Measurement (ESEM)}, 2018. DOI: 10.1145/3239235.3268920.

\bibitem{peffers2007}\dag\ K. Peffers, T. Tuunanen, M. A. Rothenberger, and S. Chatterjee, ``A design science research methodology for information systems research,'' \textit{Journal of Management Information Systems}, vol. 24, no. 3, pp. 45--78, 2007. DOI: 10.2753/MIS0742-1222240302.

\bibitem{ralph2021} P. Ralph \textit{et al.}, ``Empirical standards for software engineering research,'' arXiv:2010.03525, 2021.

\bibitem{rezaeinasab2022}\dag\ A. Rezaei Nasab, M. Shahin, S. A. Hoseyni Raviz, P. Liang, A. Mashmool, and V. Lenarduzzi, ``An empirical study of security practices for microservices systems,'' \textit{Journal of Systems and Software}, vol. 192, 2022. DOI: 10.1016/j.jss.2022.111563.

\bibitem{runeson2009}\dag\ P. Runeson and M. Höst, ``Guidelines for conducting and reporting case study research in software engineering,'' \textit{Empirical Software Engineering}, vol. 14, no. 2, pp. 131--164, 2009. DOI: 10.1007/s10664-008-9102-8.

\bibitem{vasilescu2015}\dag\ B. Vasilescu, Y. Yu, H. Wang, P. Devanbu, and V. Filkov, ``Quality and productivity outcomes relating to continuous integration in GitHub,'' in \textit{Proc. 10th Joint Meeting on Foundations of Software Engineering (ESEC/FSE)}, 2015, pp. 805--816. DOI: 10.1145/2786805.2786850.

\bibitem{wohlin2012} C. Wohlin, P. Runeson, M. Höst, M. C. Ohlsson, B. Regnell, and A. Wesslén, \textit{Experimentation in Software Engineering}. Berlin: Springer, 2012.

\bibitem{wohlin2014ease}\dag\ C. Wohlin, ``Guidelines for snowballing in systematic literature studies and a replication in software engineering,'' in \textit{Proc. 18th Int. Conf. Evaluation and Assessment in Software Engineering (EASE)}, 2014, pp. 1--10. DOI: 10.1145/2601248.2601268.

\bibitem{zampetti2020}\dag\ F. Zampetti, C. Vassallo, S. Panichella, G. Canfora, H. Gall, and M. Di Penta, ``An empirical characterization of bad practices in continuous integration,'' \textit{Empirical Software Engineering}, vol. 25, no. 2, pp. 1095--1135, 2020. DOI: 10.1007/s10664-019-09785-8.

\end{thebibliography}
